\chapter{结论}

\section{全文工作总结}

本文以同济大学 Unix V6++ 教学操作系统为对象，在团队整体 Rust 重构和 RISC-V 迁移工作的基础上，重点研究内存管理、文件系统及其相关块 I/O 路径。论文没有把所有内核模块按说明书方式逐项展开，而是围绕资源生命周期、错误传播和平台边界三个问题，分析 main、refactor/x86 和 refactor/riscv 三个分支之间的主要差异。

在内存管理模块中，Rust 版本重新组织了物理页、页级分配、内核堆和交换区管理。物理页以统一结构表示，页框号、物理地址和地址区间通过类型化抽象区分；Buddy 分配器负责页级连续空间分配，Slab 分配器负责内核小对象分配；交换区管理以显式空闲区间表组织分配、释放和合并逻辑。上述改动保留了原系统页式管理和交换机制，同时使资源粒度和释放路径更容易审查。

在文件系统模块中，Rust 版本保留 Unix V6 风格的超级块、inode 缓存、目录项查找、打开文件表和多级块映射机制，重构重点放在对象引用关系和错误处理路径上。内存 inode、文件对象和文件描述符表通过包装类型表达共享关系，资源释放依赖析构路径完成；内部错误使用 \texttt{Result<T, E>} 与 PosixError 组织，并在系统调用边界转换为符合 Unix 语义的错误返回。缓冲区缓存和块设备接口被并入文件系统路径分析，用于说明文件系统如何与具体设备驱动解耦。

在 RISC-V 迁移和测试部分，本文围绕内存管理和文件系统所依赖的底层条件，分析了 Sv39 分页、UART 串口和 VirtIO 块设备路径。测试结果表明，Rust IA32 和 Rust RISC-V 两个版本均能完成文件系统装载、shell 运行、文件读写和用户程序执行。系统调用、malloc/free 和文件 I/O 测试说明，Rust 重构没有在主路径上引入明显失控的开销；RISC-V 版本在较大文件 I/O 场景下仍主要受 VirtIO 驱动和模拟平台影响。

\section{不足与后续工作}

本文仍存在若干不足。研究边界方面，Unix V6++ 的 Rust 重构和 RISC-V 迁移是团队协作完成的系统工程，本文重点讨论内存管理、文件系统以及二者依赖的块 I/O 路径；进程管理、系统调用分发、中断处理和终端等模块只在接口依赖和测试现象层面说明。内存管理方面，当前交换区采用较直接的空闲区间管理策略，Buddy 阶数和 Slab 槽位范围主要由静态常量确定，RISC-V 版本也尚未利用大页降低页表开销。文件系统方面，当前实现仍保留固定大小目录项和有限 inode 地址表，没有扩展多挂载点、一致性恢复和更大文件支持。

后续工作可从测试验证、性能优化和教学支撑三个方向推进。测试方面，可补充页分配回收、Slab 对象复用、交换区碎片合并、inode 引用计数、目录项边界、文件截断和缓冲区缓存一致性测试。性能方面，可继续优化 VirtIO 请求合并、中断驱动 I/O 和缓存写回策略，并引入 RISC-V cycle 计数器或更高精度计时方法。教学方面，可围绕重构后的内存管理和文件系统补充代码阅读路线、实验任务和 unsafe 边界说明，使 Unix V6++ 同时保留传统 Unix 内核机制和现代系统软件实践内容。
